
\subsubsection{6.1 What the Results Mean}

\textbf{Repair is bounded by signal extraction, not information transfer.} The cleanest account of the inverse correlation is that a model repairing its own code has to identify one edit, and the feedback's job is to make that edit obvious rather than to be complete. A traceback that says a NameError occurred at line 5 has done the job in twelve words. Pyright's diagnostic tree contains the same fact somewhere, along with forty others, and locating it is now the model's problem. Information content went up and extractable signal went down. If this is right, the productive engineering target is the feedback presentation layer rather than the verifier's rigor --- and that layer is currently the least examined part of every repair loop we know of.

\textbf{The correctness-optimal and efficiency-optimal verifiers are different tools.} Execution monitoring wins on $\Delta pass$@1 by approximately $2\times$; static analysis wins on correctness-per-second by an estimated $16\times$. Neither number is more real than the other, and the choice between them is a deployment question rather than a ranking. A batch pipeline with slack latency should use execution monitoring. An interactive assistant that must return in tens of milliseconds should use static analysis and accept the lower hit rate. The field's habit of reporting correctness without cost has made this a question nobody had the numbers to ask.

\textbf{There is a model-tier floor for SMT-integrated repair.} GPT-4o-mini's 6.3\% constraint-extraction rate is not a tuning problem. Translating an informal docstring into Z3 constraints is a formal reasoning task that this tier of model does not perform reliably, and a repair pipeline that depends on it will fail on 19 of every 20 problems before the solver runs.

\subsubsection{6.2 Limitations}

\textbf{The overhead measurements come from a calibrated mock run.} An API key was unavailable when the efficiency experiment executed, so overhead was drawn from log-normal distributions parameterized by empirical priors measured in the specificity experiment and by published tool profiles. The ordering static < type < execution < SMT is robust across any realistic parameterization. The absolute ratios --- 6.64 and 0.41 --- should be read as calibrated estimates, and a live run is a straightforward replication not yet performed. As noted in Table 5.5, the $\Delta pass$@1 column values and the efficiency ratios derive from separate synthetic distributions in the mock simulation and are not arithmetically constrained to agree; the efficiency ratios are the gate metric and the primary output.

\textbf{Character count conflates precision with verbosity.} Pyright's 24,358 characters are not 24,358 characters of insight about the failing test --- they are a JSON serialization that includes unused-import notices and style warnings. Our $\rho$ = -1.0 may therefore reflect a format effect rather than a genuine relationship between formal precision and repair utility, and we cannot separate the two with this design. The practical finding survives regardless: injecting raw Pyright JSON into a repair prompt reduces repair effectiveness in this setup, whatever the mechanism.

\textbf{Truncation may be doing the work.} Four thousand characters of Pyright's output reaches the model and roughly twenty thousand do not. A naive prefix cut can discard the decisive diagnostic. We cannot currently distinguish ''verbose feedback is hard to use'' from ''we cut off the useful part,'' and these have different fixes.

\textbf{Z3's position is probably partly artifactual.} Z3 tops the repair ranking on a set where it produced real constraints for 8 of 126 problems and a fallback string otherwise. Those 8 may be structurally simpler, which would inflate its rate for reasons unrelated to feedback quality.

\textbf{One backbone.} Everything here is conditional on GPT-4o-mini's context-parsing capacity. A model that parses structured JSON well might extract Pyright's decisive diagnostic more reliably and reverse the sign, which would invert the practical recommendation rather than merely weaken it.

\textbf{The repair results are effectively HumanEval-only.} MBPP's uniform 0\% floor means the ordering in Section 5.4 rests on 46 HumanEval failures. The confidence intervals overlap, and no pairwise gap is individually significant; what carries the finding is the perfect monotonicity across four categories and its persistence under stratification.

\subsubsection{6.3 Impact Statement}

This work aims to make verifier selection in LLM repair pipelines an empirical decision rather than an intuitive one, and its immediate effect should be to reduce wasted computation --- teams currently paying an order of magnitude in latency for verbose feedback their model cannot exploit have a reason to reconsider.

Two risks deserve naming. Our efficiency framing could encourage optimizing for correctness-per-second in settings where absolute correctness is what matters; in safety-relevant code, the improvement from execution monitoring is worth its 800 milliseconds and the ratio is the wrong metric to optimize. And our results are conditional on one model tier, so treating ''avoid verbose static analysis feedback'' as a general rule would overextend them.

\subsubsection{6.4 What We Would Deploy}

Use execution monitoring when correctness per attempt is the binding constraint and latency is not. Use static analysis when throughput or interactive latency binds --- accepting roughly half the correctness gain for roughly a seventeenth of the cost. Do not inject raw Pyright JSON; summarize it first. Do not build an SMT stage on a GPT-4o-mini-class extractor. And if problems look like MBPP rather than HumanEval, put the full problem specification in the repair prompt before spending any effort choosing between verifiers --- at that point the verifier is not what is limiting you.

